% Part: first-order-logic
% Chapter: sequent-calculus
% Section: equality

\documentclass[../../../include/open-logic-section]{subfiles}

\begin{document}

\olfileid[fr]{fol}{seq}{ide}

\olsection{\usetoken{P}{derivation} avec l'\usetoken{s}{identity}}

Les !!{derivation}s avec l'!!{identity} nécessitent des séquents
initiaux et des règles d'inférence supplémentaires.

\begin{defn}[Séquents initiaux pour $\eq$]
Si $t$ est un terme clos, alors ${} \Sequent \eq[t][t]$ est un
séquent initial.
\end{defn}

Les règles pour $\eq$ sont les suivantes, où $t_1$ et $t_2$ sont
des termes clos :

\begin{defish}
\Axiom$ \eq[t_1][t_2], \Gamma \fCenter \Delta, !A(t_1) $
\RightLabel{$\eq$}
\UnaryInf$\eq[t_1][t_2], \Gamma \fCenter \Delta, !A(t_2)$
\DisplayProof
\hfill
\Axiom$\eq[t_1][t_2], \Gamma \fCenter \Delta, !A(t_2) $
\RightLabel{$\eq$}
\UnaryInf$\eq[t_1][t_2], \Gamma  \fCenter \Delta, !A(t_1)$
\DisplayProof
\end{defish}

\begin{ex}
Si $s$ et $t$ sont des termes clos, alors
$\eq[s][t], !A(s) \Proves !A(t)$ :
\begin{prooftree}
\Axiom$ !A(s) \fCenter !A(s)$
\RightLabel{\LeftR{\Weakening}}
\UnaryInf$\eq[s][t], !A(s)  \fCenter !A(s)$
\RightLabel{$\eq$}
\UnaryInf$\eq[s][t], !A(s)  \fCenter !A(t)$
\end{prooftree}
On reconnaît le principe de substituabilité des égaux, également
appelé loi de Leibniz.

Dans $\Log{LK}$, on démontre que $\eq$ est symétrique et transitive :
\begin{prooftree}
\Axiom$ \fCenter \eq[t_1][t_1] $
\RightLabel{\LeftR{\Weakening}}
\UnaryInf$ \eq[t_1][t_2] \fCenter \eq[t_1][t_1] $
\RightLabel{$\eq$}
\UnaryInf$ \eq[t_1][t_2] \fCenter \eq[t_2][t_1]$
\DisplayProof\qquad\bottomAlignProof
\Axiom$ \eq[t_1][t_2] \fCenter \eq[t_1][t_2] $
\RightLabel{\LeftR{\Weakening}}
\UnaryInf$\eq[t_2][t_3], \eq[t_1][t_2]  \fCenter \eq[t_1][t_2] $
\RightLabel{$\eq$}
\UnaryInf$\eq[t_2][t_3], \eq[t_1][t_2]  \fCenter \eq[t_1][t_3]$
\RightLabel{\LeftR{\Exchange}}
\UnaryInf$\eq[t_1][t_2], \eq[t_2][t_3]  \fCenter \eq[t_1][t_3]$
\end{prooftree}
Dans la !!{derivation} de gauche, la !!{formula}~$\eq[x][t_1]$
joue le rôle de $!A(x)$. À droite, nous prenons
$!A(x)$ égal à $\eq[t_1][x]$.
\end{ex}

\begin{prob}
Construire des !!{derivation}s des séquents suivants :
\begin{enumerate}
\item $\Sequent \lforall[x][\lforall[y][((x = y \land !A(x)) \lif !A(y))]]$
\item $\lexists[x][!A(x)] \land \lforall[y][\lforall[z][((!A(y) \land
    !A(z)) \lif y = z)]] \Sequent 
\lexists[x][(!A(x) \land \lforall[y][(!A(y) \lif y = x)])]$
\end{enumerate}
\end{prob}

\end{document}
